\begin{center}
	\large{\textbf{COEFFICIENT MATRIX FOR TWO AND THREE QUBIT STATES ENTANGLEMENT MEASURE}}
\end{center}
%\begin{flushleft}
\begin{minipage}{4cm}
	\begin{align*}
		&\text{Name} \qquad &&: \text{Tamara Pingki} \\
		&\text{NRP} \qquad &&: \text{6001241001} \\
		&\text{Department} \qquad &&: \text{Fisika} \\
		&\text{Supervisor} \qquad &&: \text{Prof. Agus Purwanto, D.Sc.} \\
	\end{align*}
\end{minipage} \\
\begin{center}
	\large\textbf{ABSTRACT} \\  
\end{center}
Quantum entanglement (\textit{quantum entanglement}) is a fundamental resource in quantum information processing. Although the classification of entanglement in bipartite systems is well understood, identifying entanglement classes in multipartite systems remains a complex problem due to the increasing number of entanglement configurations as the number of qubits increases. This study aims to formulate a coefficient matrix-based theoretical framework to identify and classify entanglement in two- and three-qubit systems. The research method is carried out analytically through the formulation of a coefficient matrix constructed from the amplitudes of quantum states. For a two-qubit system, the coefficient matrix formulation is $\mathrm{Tr},\rho \tilde{\rho} = 4 \det \rho_A$, while for a three-qubit system, the relationship is $\mathrm{Tr}\,\rho_{AB} \tilde{\rho}_{AB}= 4 |\det C_{AB0}|^2 + 4 |\det C_{AB1}|^2 + 2 (\det C_{AB0}^1 + \det C_{AB1}^0)(\det C^{*1}_{AB0} + \det C^{*0}_{AB1})$. This formulation is then applied to various two- and three-qubit quantum states to determine their entanglement distributions. The results of this study show that the coefficient matrix-based approach can be used as a simpler and more efficient alternative method to characterize entanglement classes in two- and three-qubit systems.

\noindent
\begin{minipage}{1cm}
	\begin{align*}
		\text{\textbf{Kata kunci}}:\hspace{5mm}&\text{Coefficient Matrix, Entanglement,Two-Qubit, Three-Qubit}% \\ &\text{} \\&\text{}
	\end{align*}
\end{minipage}


\newpage
\thispagestyle{empty}
\begin{center}
	\topskip0pt
	\vspace*{\fill}
	\textit{\textbf{"Halaman ini sengaja dikosongkan"}}
	\vspace*{\fill}
\end{center}
\pagebreak